\begin{quote}
Pri vývoji prekladača sme využívali nástroje umelej inteligencie ako prostriedok na zefektívnenie práce a pochopenie problematiky. Najčastejšie sme pracovali s modelom Gemini, no pri riešení čiastkových problémov sme využívali aj nástroje Claude a ChatGPT. S umelou inteligenciou sme diskutovali o správnom pochopení častí zadania a špecifikácie jazyka IFJ26. V situáciách, keď nám niektoré koncepty alebo požiadavky zadania neboli jasné, alebo sme nevedeli ako konkrétnu časť riešiť, prípadne sme ju nevedeli nájsť v zadaní, nám AI dokázala zrozumiteľne vysvetliť podstatu problému a pomohla nám zorientovať sa v zadaní. Okrem konzultácií k návrhu sme tieto nástroje zapájali aj pri ladení a debugovaní kódu a pri vytváraní testov na overovanie správneho aj chybového správania prekladača. Práve v prípade testov nám využitie ai prišlo ako najviac praktické a dokázali sme efektívne debuovať kód a spojiť časti, ktoré sme robili oddelene dokopy.

\end{quote}

